%%=============================================================================
%% Voorwoord
%%=============================================================================

\chapter*{\IfLanguageName{dutch}{Woord vooraf}{Preface}}%
\label{ch:voorwoord}

%% TODO:
%% Het voorwoord is het enige deel van de bachelorproef waar je vanuit je
%% eigen standpunt (``ik-vorm'') mag schrijven. Je kan hier bv. motiveren
%% waarom jij het onderwerp wil bespreken.
%% Vergeet ook niet te bedanken wie je geholpen/gesteund/... heeft
De keuze voor mijn onderwerp was voor mij geen toeval. 
De gezondheidszorg is een domein dat mij nauw aan het hart ligt. 
Verschillende gezins- en familieleden zijn actief in de zorgsector, waardoor ik van jongs af aan in aanraking ben gekomen met de uitdagingen en de meerwaarde van kwalitatieve zorgverlening. 
Daarnaast liep ik stage bij de IT-dienst van een ziekenhuis, een ervaring die mij toonde hoe groot de impact van goed ontworpen digitale tools kan zijn op de dagelijkse werking van zorgprofessionals. 
De combinatie van deze persoonlijke achtergrond en mijn interesse in toegepaste informatica maakte dit onderzoeksonderwerp voor mij dan ook een evidente en gemotiveerde keuze.
\medskip

Deze bachelorproef was niet mogelijk geweest zonder de steun en medewerking van een aantal mensen die ik graag uitdrukkelijk wil bedanken.\medskip

In de eerste plaats gaat mijn oprechte dank uit naar mijn promotor, dhr. Benjamin Vertonghen, voor zijn begeleiding doorheen het volledige traject. 
Zijn constructieve feedback en bereikbaarheid hebben mij geholpen om dit onderzoek stap voor stap vorm te geven.
Ook mijn co-promotor, mevr. S. Van Rampelberg, verdient een woord van dank. 
Haar inhoudelijke expertise rond de toegankelijkheid van de eerstelijnszorg was een waardevolle bijdrage in dit onderzoek. 
Tien dagen voor de finale deadline van indiening ontving ik echter een mail met de mededeling dat zij niet langer deel uitmaakte van het project en bijgevolg ook niet aanwezig zal zijn op de verdediging van mijn bachelorproef. 
Dit was een onverwachte en niet aangename mededeling, maar ik waardeer de samenwerking die we hebben gehad.\medskip 

Verder wil ik alle deelnemers aan de gebruikerstesten hartelijk bedanken. 
Wetende dat huisartsen en praktijkmedewerkers  een drukke agenda hebben, namen zij toch de tijd om mijn applicatie te testen en doordachte feedback te geven. 
Zonder hun inbreng was een evaluatie van de proof-of-concept interface niet mogelijk geweest.
Ten slotte bedank ik ook de stuurgroep voor hun feedback op het ontwerp en de inhoud van de interface. 
Hun betrokkenheid heeft dit onderzoek mede mogelijk gemaakt.
